% !TeX root = ../../Book3.tex
\section{Blaschke 乘积与零点}
\label{b3:sec:2803}

上一节把 \(H^2(\mathbb D)\) 中的解析函数与圆周上的平方可积边界函数联系起来。现在要问：内部的零点怎样影响这份边界资料？直接除以 \(z-a\) 虽能消去一个零点，却会改变圆周各处的模。Blaschke 因子在单位圆周模为一，因而适合把零点从函数中分离出来，同时保留边界模和 Hardy 范数。

\ProseParagraphBreak

\chapref{b3:ch:18}已经证明有界解析函数的零点满足 Blaschke 条件，并从满足该条件的零点列构造了乘积。本节不再重做那套构造，而要补足两件事：平方平均有界为何也给出同一零点条件；无穷乘积为何仍有模为一的边界值，从而可以在 \(H^2\) 中作等范数除法。仍取 \(\mathbb D=\{\,z\in\mathbb C\mid |z|<1\,\}\)，圆周测度为 \(\mathrm dm=\mathrm d\theta/(2\pi)\)，\(f^*\) 表示\cref{b3:thm:28-hardy-boundary}建立的 \(L^2\) 边界函数。

\subsection{零点列与 Blaschke 条件}
\label{b3:subsec:280301}

解析性只禁止零点在圆盘内部聚集；它并未限制零点沿边界聚集得多快。\(H^2\) 的额外约束来自所有内部圆周平方平均的共同上界。Jensen 公式恰好把这个平均界转换为零点的加权总量。

\Needspace{6\baselineskip}
\begin{theorem}[Hardy 函数的零点条件]\label{b3:thm:28-hardy-zero-condition}
设 \(f\in H^2(\mathbb D)\)、\(f\not\equiv0\)，原点是它的 \(m\) 重零点，其中 \(m\ge0\)，并记第一个非零 Taylor 系数为 \(c_m=f^{(m)}(0)/m!\)。把全部非零零点按重数列为 \((a_j)\)，允许有限列或空列。则
\begin{equation}\label{b3:eq:28-zero-budget}
 \sum_j-\log|a_j|
 \le \log\frac{\|f\|_{H^2}}{|c_m|}<\infty,
 \quad
 \sum_j(1-|a_j|)<\infty.
\end{equation}
\end{theorem}
\begin{proof}
固定 \(0<r<1\)。函数 \(f\) 在闭盘 \(|z|\le r\) 的开邻域全纯，所以\cref{b3:thm:18-jensen}适用；即使观察圆周上有零点，其对数模仍可积。该公式给出
\[
 \int\log|f(re^{\mathrm i\theta})|\,\mathrm dm
 =\log|c_m|+m\log r
       +\sum_{|a_j|<r}\log\frac r{|a_j|}.
\]
对数的凹性以及 \(m(\partial\mathbb D)=1\) 给出
\[
 \int\log|f(re^{\mathrm i\theta})|\,\mathrm dm
 \le\frac12\log\int|f(re^{\mathrm i\theta})|^2\,\mathrm dm
 \le\log\|f\|_{H^2}.
\]
这是\eqref{b3:eq:18-mean-power-bound}在指数二时的版本。故
\[
 \sum_j\left(\log\frac r{|a_j|}\right)_+
 \le\log\frac{\|f\|_{H^2}}{|c_m|}-m\log r.
\]
左侧每一项随 \(r\uparrow1\) 非减，趋于 \(-\log|a_j|\)。对计数测度使用单调收敛，且注意 \(m\log r\to0\)，便得到第一个不等式。最后利用 \(1-t\le-\log t\)（\(0<t<1\)）得到 Blaschke 条件。
\end{proof}

原点零点由有限整数 \(m\) 单独记录，不能把 \(-\log0\) 塞入和式。式\eqref{b3:eq:28-zero-budget}还保留了定量信息：若 \(f(0)\ne0\)，则
\[
 \prod_j|a_j|\ge\frac{|f(0)|}{\|f\|_{H^2}}>0.
\]
中心值相对范数越大，零点可以消耗的对数总量就越少。例如 \(|f(0)|=\|f\|_{H^2}\) 时，系数等距已经迫使所有正次系数为零，函数是非零常数；零点预算为零与这个结论一致。

\ProseParagraphBreak

求和必须计入重数。一个点即使非常靠近边界，若把它重复足够多次，仍可使总和发散。这里控制的并非圆盘中的几何点集，而是解析函数所携带的计重数零点列。

\subsection{Blaschke 乘积的收敛}
\label{b3:subsec:280302}

先恢复乘积的归一化。对 \(a\in\mathbb D\setminus\{0\}\)，取
\begin{equation}\label{b3:eq:28-blaschke-factor}
 b_a(z)=\frac{|a|}{a}\frac{a-z}{1-\overline a z},
 \quad b_a(0)=|a|.
\end{equation}
它与\cref{b3:def:06-finite-blaschke}使用的 \((z-a)/(1-\overline a z)\) 只差一个模一常数，所以零点与模都相同。有限乘积中这个常数可以并入整体相位；无穷乘积却必须固定相位，使靠近边界的因子在内部紧集上趋于一。

\ProseParagraphBreak

具体说，设 \((a_j)\subset\mathbb D\setminus\{0\}\) 按有限重数重复，满足 \(\sum_j(1-|a_j|)<\infty\)，另取整数 \(m\ge0\) 和 \(|\eta|=1\)。这个条件本身已保证任何内部闭盘中只有有限项，否则每项缺陷都有共同正下界。\cref{b3:thm:18-blaschke-zeros}因而适用，给出
\begin{equation}\label{b3:eq:28-blaschke-product}
 B(z)=\eta z^m\prod_j b_{a_j}(z).
\end{equation}
乘积在圆盘内局部一致收敛，不恒为零，其零点恰为所列零点，且 \(|B(z)|\le1\)。空乘积取一。回忆收敛的关键估计：若 \(a=\rho e^{\mathrm i\varphi}\)，则
\[
 \begin{aligned}
 1-b_a(z)&=(1-\rho)\frac{1+e^{-\mathrm i\varphi}z}
                              {1-\rho e^{-\mathrm i\varphi}z},\\
 \sup_{|z|\le s}|1-b_a(z)|&\le\frac{1+s}{1-s}(1-|a|)
 \quad(0<s<1).
 \end{aligned}
\]
由此得到尾因子的局部对数级数绝对一致收敛；它的指数无零点，有限头部则保留指定的重数。这个论证也说明不能去掉式\eqref{b3:eq:28-blaschke-factor}中的相位，只凭每个因子模不超过一来断言乘积收敛。

\ProseParagraphBreak

有限乘积在每个边界点模为一，但局部一致极限只控制圆盘内部，不能直接传递这项边界性质。需要结合刚刚建立的 \(L^2\) 边界理论。为此先给出所用的术语。

\Needspace{6\baselineskip}
\begin{definition}\label{b3:def:28-inner}
设 \(I\) 在 \(\mathbb D\) 全纯且有界。若它的 \(L^2(\partial\mathbb D,m)\) 边界函数满足 \(|I^*|=1\) 几乎处处，则称 \(I\) 为\term[nei-hanshu]{内函数}[inner function]。这里的边界函数由 \(I\in H^2(\mathbb D)\) 及径向 \(L^2\) 收敛确定，不要求连续到闭圆盘。
\end{definition}

\Needspace{6\baselineskip}
\begin{theorem}[Blaschke 乘积的内函数性质]\label{b3:thm:28-blaschke-inner}
设 \(m\ge0\) 为整数，\(|\eta|=1\)，且非零点列 \((a_j)\subset\mathbb D\) 按有限重数重复并满足 Blaschke 条件。式\eqref{b3:eq:28-blaschke-product}定义的有限或无穷乘积 \(B\) 是内函数，且
\[
 \|B\|_{H^2}=1,
 \quad
 \lim_{r\uparrow1}\int-\log|B(re^{\mathrm i\theta})|\,\mathrm dm=0.
\]
\end{theorem}
\begin{proof}
由构造，\(B\) 在盘内模不超过一，故属于 \(H^2\)，并存在 \(L^2\) 边界函数 \(B^*\)。Blaschke 条件还给出 \(\sum_j-\log|a_j|<\infty\)：除了有限项以外 \(|a_j|\ge1/2\)，而在 \(1/2\le t<1\) 时有 \(-\log t\le2(1-t)\)。因此 \(B\) 的第一个非零 Taylor 系数为
\[
 c_m=\eta\prod_j|a_j|\ne0.
\]
对每个 \(0<r<1\) 使用 Jensen 公式，得到
\[
 \int\log|B(re^{\mathrm i\theta})|\,\mathrm dm
 =m\log r+\sum_j\log|a_j|
             +\sum_{|a_j|<r}\log\frac r{|a_j|}.
\]
最后一个非负和随 \(r\uparrow1\) 单调趋于 \(\sum_j-\log|a_j|\)。两个有限级数恰好抵消，\(m\log r\to0\)，所以对数平均趋于零。

由径向 \(L^2\) 收敛，可以选 \(r_k\uparrow1\)，使 \(B(r_ke^{\mathrm i\theta})\to B^*(e^{\mathrm i\theta})\) 几乎处处。因内部模不超过一，\(|B^*|\le1\) 几乎处处。对非负函数 \(-\log|B(r_ke^{\mathrm i\theta})|\) 使用 Fatou 引理；在零值处规定 \(-\log0=+\infty\)，得到
\[
 0\le\int-\log|B^*|\,\mathrm dm
 \le\liminf_{k\to\infty}\int-\log|B(r_ke^{\mathrm i\theta})|\,\mathrm dm=0.
\]
故 \(-\log|B^*|=0\) 几乎处处，即 \(|B^*|=1\)。边界等距于是给 \(\|B\|_{H^2}^2=\int|B^*|^2\,\mathrm dm=1\)。
\end{proof}

这里没有把一个无控制的对数极限移进积分。我们先由 Jensen 公式算出非负积分趋于零，再用 Fatou 引理判断极限。所需的逐点收敛也只有 \(L^2\) 收敛所提供的一条子列；无需预先假设所有径向极限都存在。

\ProseParagraphBreak

乘积不仅在内部紧集收敛，在 Hardy 范数下也收敛。设零点列无限，取
\[
 B_N=\eta z^m\prod_{j=1}^N b_{a_j},
 \quad T_N=\prod_{j>N}b_{a_j},
 \quad B=B_NT_N.
\]
两个乘积的边界模均为一。利用\cref{b3:prop:28-hardy-multiplier}的边界乘法公式，以及内积对第一变量线性的约定，有
\[
 \langle B,B_N\rangle_{H^2}
 =\int B_N^*T_N^*\overline{B_N^*}\,\mathrm dm
 =\int T_N^*\,\mathrm dm=T_N(0)=\prod_{j>N}|a_j|.
\]
最后一步是边界 Fourier 零次系数等于解析函数的常数项。因此
\begin{equation}\label{b3:eq:28-blaschke-hardy-convergence}
 \|B-B_N\|_{H^2}^2
 =2\left(1-\prod_{j>N}|a_j|\right)\longrightarrow0.
\end{equation}
尾部对数和趋于零，正是尾乘积在原点趋于一的原因。这项平均意义的逼近不声称在整个闭圆盘上一致收敛。

\subsection{零点分解与 Hardy 范数}
\label{b3:subsec:280303}

对有界函数，第六章和第十八章已用最大模原理证明除以 Blaschke 因子仍有相同模界。\(H^2\) 函数可能无界，此时要控制的是圆周积分。先除有限个零点，再让这批零点逐渐覆盖全部零点，可以避免在尚未证明可积性时直接对无穷商取边界值。

\Needspace{6\baselineskip}
\begin{lemma}[有限零点的等范数除法]\label{b3:lem:28-finite-blaschke-division}
设 \(f\in H^2(\mathbb D)\)，\(P\) 为有限 Blaschke 乘积，且 \(P\) 的每个零点在 \(f\) 中至少具有相同重数。则可去延拓后的商 \(q=f/P\) 属于 \(H^2(\mathbb D)\)，并且 \(\|q\|_{H^2}=\|f\|_{H^2}\)。空乘积或模一常数也允许。
\end{lemma}
\begin{proof}
零点的重数条件使所有商奇点可去，故 \(q\) 在整个圆盘全纯。取 \(r\) 大于 \(P\) 的全部零点的模，令
\[
 \mu_r=\min_{|z|=r}|P(z)|>0.
\]
有限乘积连续到闭圆盘且在单位圆周模为一，因此 \(\mu_r\to1\)。圆周上有
\[
 \int|q(re^{\mathrm i\theta})|^2\,\mathrm dm
 \le\mu_r^{-2}\int|f(re^{\mathrm i\theta})|^2\,\mathrm dm
 \le\mu_r^{-2}\|f\|_{H^2}^2.
\]
解析函数的平方平均由其 Taylor 系数公式可知随半径非减；令 \(r\uparrow1\)，便得到 \(q\in H^2\) 及 \(\|q\|_{H^2}\le\|f\|_{H^2}\)。反过来，由 \(f=Pq\) 和 \(|P|\le1\)，每个内部圆周上的平方平均都满足反向范数所需的不等式，故 \(\|f\|_{H^2}\le\|q\|_{H^2}\)。
\end{proof}

\Needspace{6\baselineskip}
\begin{theorem}[Hardy 函数的零点分解]\label{b3:thm:28-zero-factorization}
设 \(f\in H^2(\mathbb D)\)、\(f\not\equiv0\)。令 \(B\) 为按 \(f\) 的全部零点及重数构造的 Blaschke 乘积，原点幂与非零因子的归一化如式\eqref{b3:eq:28-blaschke-product}，整体相位任选。则
\[
 f=Bg,
 \quad g\in H^2(\mathbb D),
 \quad g(z)\ne0\quad(z\in\mathbb D),
\]
并且
\begin{equation}\label{b3:eq:28-zero-factorization-boundary}
 \begin{gathered}
 f^*=B^*g^*\quad\text{几乎处处},
 \quad |f^*|=|g^*|\quad\text{几乎处处},\\
 \|f\|_{H^2}=\|g\|_{H^2}.
 \end{gathered}
\end{equation}
\end{theorem}
\begin{proof}
\cref{b3:thm:28-hardy-zero-condition}保证相应 Blaschke 乘积存在。由于 \(B\) 与 \(f\) 具有完全相同的零点及重数，\(g=f/B\) 的商奇点全部可去，补值后在圆盘内全纯且无零点。关键尚在于证明它属于 \(H^2\)。

零点有限时，上一引理已给出结论。设非零零点列无限，并令 \(B_N\) 为包括全部原点幂的有限部分乘积。\(g_N=f/B_N\) 全纯，且上一引理给出 \(\|g_N\|_{H^2}=\|f\|_{H^2}\)。下面说明 \(g_N\to g\) 局部一致，包括共同零点处。

固定内部闭盘 \(|z|\le s\)，选 \(s<R<1\)，使圆周 \(|z|=R\) 上没有 \(B\) 的零点；零点离散，因而可避开它们的模。此圆周上的 \(|B|\) 有正下界，而 \(B_N\to B\) 一致，所以 \(g_N=f/B_N\to f/B=g\) 在该圆周上一致。最大模原理用于全纯差 \(g_N-g\)，将收敛传到 \(|z|\le R\)。于是对每个 \(0<s<1\)，
\[
 \int|g(se^{\mathrm i\theta})|^2\,\mathrm dm
 =\lim_{N\to\infty}\int|g_N(se^{\mathrm i\theta})|^2\,\mathrm dm
 \le\|f\|_{H^2}^2.
\]
这就证明 \(g\in H^2\) 及 \(\|g\|_{H^2}\le\|f\|_{H^2}\)。由 \(f=Bg\)、\(|B|\le1\) 又得反向不等式。

现在 \(B\) 有界解析且 \(g\in H^2\)，\cref{b3:prop:28-hardy-multiplier}给出 \(f^*=B^*g^*\)。\cref{b3:thm:28-blaschke-inner}给出 \(|B^*|=1\) 几乎处处，遂得到边界模相等及式\eqref{b3:eq:28-zero-factorization-boundary}。
\end{proof}

这个结论把两种作用分开：\(B\) 记录内部零点，\(g\) 保留全部边界模；但 \(B^*\) 的相位仍会改变边界函数本身。范数相等不是说函数值相同，也不是说 \(g\) 有界。整体相位的自由只会把 \(g\) 乘上一个相反的模一常数，不影响这些结论。

\subsection{零点约束与边界增长}
\label{b3:subsec:280304}

本节的必要条件也是零点列可由某个非零 \(H^2\) 函数实现的充分条件：满足 Blaschke 条件时，所构造的 \(B\) 本身就属于 \(H^2\)。因此，从“有界解析函数”放宽到“平方平均有界的解析函数”，并没有增加可实现的零点列；增加的是同一零点列可以搭配的无零点余因子。

\ProseParagraphBreak

例如取 \(\alpha>0\) 及 \(a_n=1-n^{-\alpha}\)、\(n\ge2\)。这些点能成为某个非零 \(H^2\) 函数的全部零点，当且仅当
\[
 \sum_{n=2}^{\infty}(1-a_n)=\sum_{n=2}^{\infty}n^{-\alpha}<\infty,
 \quad\text{即}\quad\alpha>1.
\]
指数越大，同样多的零点越靠近边界，每个零点在 Jensen 公式中付出的权重越小。边界聚集本身并不被禁止；受限的是带重数的总距离缺陷。

\Needspace{6\baselineskip}
\begin{corollary}[内部零点数的边界尺度]\label{b3:cor:28-zero-count-scale}
设 \(f\in H^2(\mathbb D)\)、\(f\not\equiv0\)，令 \(n_f(r)\) 为 \(|z|\le r\) 中的计重数零点个数，\(0<r<1\)。则
\[
 (1-r)\cdot n_f(r)\longrightarrow0\quad(r\uparrow1).
\]
\end{corollary}
\begin{proof}
原点的有限重数乘以 \(1-r\) 后趋于零，只需处理非零零点 \((a_j)\)。给定 \(\varepsilon>0\)，由 Blaschke 条件选 \(N\)，使 \(\sum_{j>N}(1-|a_j|)<\varepsilon\)。对 \(|a_j|\le r\) 有 \(1-r\le1-|a_j|\)，故
\[
 (1-r)\#\{\,j:|a_j|\le r\,\}
 \le N(1-r)+\sum_{j>N}(1-|a_j|)<N(1-r)+\varepsilon.
\]
先令 \(r\uparrow1\)，再令 \(\varepsilon\downarrow0\)，即得结论。
\end{proof}

与之相伴，\cref{b3:thm:28-hardy-coefficients}的 Cauchy--Schwarz 估计给出
\[
 |f(z)|\le\frac{\|f\|_{H^2}}{\sqrt{1-|z|^2}}.
\]
它是单点增长的共同上界，Blaschke 条件则是全部零点的总量约束。两者承担的任务不同。乘以任意 Blaschke 乘积可以增加许多零点，同时保持 Hardy 范数与几乎处处边界模；所以只观察边界模，不能恢复零点列。

\ProseParagraphBreak

反过来，除去全部零点后也不能认为所有边界结构都消失了。回到\secref{b3:sec:1802}中见过的无零点函数，并引入正参数 \(\tau\)：
\begin{equation}\label{b3:eq:28-singular-inner-example}
 S_\tau(z)=\exp\left(-\tau\frac{1+z}{1-z}\right),
 \quad\tau>0.
\end{equation}
由于
\[
 \operatorname{Re}\frac{1+z}{1-z}
 =\frac{1-|z|^2}{|1-z|^2}>0,
\]
它在圆盘内全纯、无零点且 \(|S_\tau|\le1\)，因此属于 \(H^2\)。对 \(e^{\mathrm i\theta}\ne1\)，公式直接连续延拓到这个边界点，并有
\[
 \frac{1+e^{\mathrm i\theta}}{1-e^{\mathrm i\theta}}
 =\mathrm i\cot\frac\theta2,
 \quad
 S_\tau(re^{\mathrm i\theta})\longrightarrow
       \exp\left(-\mathrm i\tau\cot\frac\theta2\right).
\]
极限模为一。沿径向 \(L^2\) 收敛的一条几乎处处收敛子列识别边界函数，便知 \(|S_\tau^*|=1\) 几乎处处。所以 \(S_\tau\) 是一个非恒定内函数，\(\|S_\tau\|_{H^2}=1\)；但 \(S_\tau(0)=e^{-\tau}\)，且沿实轴
\[
 S_\tau(r)=\exp\left(-\tau\frac{1+r}{1-r}\right)\longrightarrow0.
\]
一个边界点上的这种行为不影响 \(L^2\) 边界范数，却使连续到闭盘的直觉失效。

\ProseParagraphBreak

这个无零点内函数的边界奇性可以直接从对数模看见：
\[
 -\log|S_\tau(re^{\mathrm i\theta})|
 =\tau\frac{1-r^2}{|1-re^{\mathrm i\theta}|^2},
 \quad
 \int-\log|S_\tau(re^{\mathrm i\theta})|\,\mathrm dm=\tau.
\]
右端是集中于边界点 \(1\) 的质量 \(\tau\) 的 Poisson 积分。对每个避开 \(1\) 的闭圆弧，核随 \(r\uparrow1\) 一致趋于零，而它的圆周积分始终为一；负对数的质量于是集中到该点。虽然 \(-\log|S_\tau^*|=0\) 几乎处处，其内部圆周积分并不趋于零。这也解释了为什么 Blaschke 内性证明中必须先用 Jensen 公式算出积分趋零，不能仅凭几乎处处边界收敛换序。

\ProseParagraphBreak

对 \(S_\tau\) 作\cref{b3:thm:28-zero-factorization}的分解时，零点乘积可取 \(B=1\)，余因子仍是 \(S_\tau\)。更一般地，\(BS_\tau\) 与 \(B\) 有相同的全部零点和相同的边界模，却是不同的解析函数。因此本节证明的是完整的“零点分离”，并未把无零点余因子的边界结构全部分类。一般的内外因子分解还须进一步区分这种集中在零测集上的边界资料；本章后续结论不以那套分类作为未证前提。
